\documentclass[10pt,landscape]{article}
\usepackage[letterpaper,landscape,margin=0.16in]{geometry}
\usepackage{multicol}
\usepackage{enumitem}
\usepackage{titlesec}
\usepackage{microtype}
\usepackage[T1]{fontenc}
\usepackage[utf8]{inputenc}
\usepackage{lmodern}
\usepackage{amsmath,amssymb}
\usepackage{ragged2e}
\pagestyle{empty}
\setlength{\parindent}{0pt}
\setlength{\parskip}{0pt}
\setlength{\columnsep}{0.12in}
\setlength{\columnseprule}{0.25pt}
\setlist[itemize]{leftmargin=0.65em,nosep,topsep=0pt,partopsep=0pt,itemsep=0pt,parsep=0pt}
\titlespacing*{\section}{0pt}{1pt}{0.5pt}
\titlespacing*{\subsection}{0pt}{0.5pt}{0pt}
\titleformat{\section}{\bfseries\fontsize{6.2}{6.2}\selectfont}{}{0em}{}
\titleformat{\subsection}{\bfseries\fontsize{5.5}{5.6}\selectfont}{}{0em}{}
\newcommand{\kw}[1]{\textbf{#1}}
\newcommand{\sep}{\vspace{0.8pt}\hrule\vspace{0.8pt}}
\newcommand{\tinytext}{\fontsize{5.75}{5.95}\selectfont}
\begin{document}
\tinytext
\RaggedRight

% ===================== PAGE 1 =====================
\begin{center}\bfseries\fontsize{7}{7}\selectfont CAS AN101 FINAL CHEAT SHEET - FRONT: WEEKS 1--6 / FOUNDATIONS, METHOD, CULTURE, KINSHIP, STATE\end{center}\vspace{-4pt}
\begin{multicols}{3}

\section*{W1: ANTHRO FOUNDATIONS}
\kw{Anthropology}: comparative study of human life/culture; makes strange familiar + familiar strange. \kw{Culture}: learned/shared/contested systems of meaning/practice; everyone has culture; not bounded like a sealed box. \kw{Social norms}: unspoken rules; violations reveal expectations. \kw{Othering}: producing ``us/them'' by language/images; makes unequal hierarchy seem natural. Watch terms: ``primitive,'' ``savage,'' ``magic,'' ``rites,'' ``tradition,'' ``tribe,'' ``natives.'' \kw{Early model}: Western/intrepid white male scholar, Western institutions, field = non-West, knowledge extraction; linked to colonialism. \kw{Noble savage}: romanticizes non-West as pure/natural but still objectifies. \kw{Great Man theory}: history via heroic individuals; anthro counters w/ structures, culture, social life. \kw{Individual choice} alone insufficient: choices shaped by norms, habitus, institutions, power. \kw{Freudian/psychology-only explanations}: anthro adds cultural/social context.
\sep
\section*{Articles always know: Miner}
\kw{Miner, Nacirema}: satire of Americans (Nacirema=American). Bathroom=shrine; medicine cabinet=charm box; dentist=holy-mouth-man; hospital=latipso. Purpose: exposes ethnocentric language; familiar practices look bizarre when described as exotic rituals. Takeaway: do not confuse descriptive distance with objectivity; language frames moral judgment. \kw{Test use}: If prompt mentions ritual/body/magic/health/othering, cite Nacirema to show relativism + critique of ethnographic voice.

\section*{W2: BOAS / RELATIVISM / RACISM}
\kw{Franz Boas}: anti-evolutionist, anti-racist; cultures must be understood historically + on their own terms. \kw{Cultural relativism}: interpret beliefs/practices in local context; not the same as approving everything. \kw{Ethnocentrism}: judging others by one's own standards; often invisible. \kw{Biological determinism}: bodies/genes determine ability, disease, criminality, intelligence, jobs. \kw{Scientific racism}: fake science ranking races; used to justify slavery/colonialism/segregation. \kw{Social evolutionism}: societies ranked from ``primitive'' to ``civilized''; abandoned because teleological, racist, empirically weak. \kw{Salvage ethnography}: urgent recording of cultures assumed to be disappearing; moral impulse but problematic (freezes cultures; ignores colonial causes; treats people as vanishing objects). \kw{Colonialism + anthro}: discipline emerged in imperial context; knowledge often helped rule. \kw{Margaret Mead}: study other societies to answer questions about one's own society.
\sep
\section*{W2: MALINOWSKI / METHOD}
\kw{Malinowski}: long-term fieldwork, learn language, live among people, document everyday life. \kw{Imponderabilia of everyday life}: small routines, gestures, rhythms; often more revealing than formal statements. \kw{Participant observation}: participate in daily life while observing social behavior; balances emic intimacy + etic analysis. \kw{Fieldnotes}: detailed records of observation, interviews, context, reflexive thoughts. \kw{Ethnography}: both method + written product. \kw{Interviews}: structured=pre-set; semi-structured=guided but flexible; unstructured=open. \kw{Key informants/interlocutors}: people who explain local meanings; not neutral representatives of whole culture. \kw{Armchair ethnography}: theorizing from others' reports; no fieldwork. \kw{Veranda ethnography}: researcher stays at colonial outpost, locals come to them; distance from everyday life. \kw{Headhunting}: example used to think beyond sensational surface to social/political meaning. \kw{Emic}: insider categories/meanings. \kw{Etic}: analytic outsider categories. Best work often combines both.
\sep
\section*{Boas Reading: Eskimo}
\kw{A Year Among the Eskimo}: field experience, learning by immersion; challenges exotic stereotypes; shows environmental adaptation, social life, stories, practical skill. Use for cultural relativism + critique of judging Arctic practices by outside comfort/standards.

\section*{W3: POSITIONALITY / HURSTON}
\kw{Insiderness/outsiderness}: not binary; degrees depending on race, class, gender, language, nationality, education, situation. \kw{Rosaldo}: grief/position can deepen understanding; emotions shape knowledge. \kw{Narayan}: ``native anthropologist'' is oversimplified; everyone has shifting identities. \kw{Hurston, Mules and Men}: Black woman anthropologist; student of Boas; Eatonville, FL; documents folklore, jokes, speech, everyday life. \kw{Dialect}: Hurston records authentic speech; controversy: some Harlem Renaissance critics saw it as stereotypes, but it preserves linguistic/cultural richness. \kw{Burden of representation}: minority author/scholar is pressured to represent whole group positively/correctly. \kw{Terms shift}: ``native''/``informant'' increasingly replaced by locals, interlocutors, participants, collaborators.
\sep
\section*{W3: SAID / NADER / KNOWLEDGE}
\kw{Said, Orientalism}: Western knowledge system constructs ``the East/Orient'' as irrational, sensual, despotic, backward; West becomes rational, modern, superior. Knowledge is not neutral; it authorizes rule/intervention. \kw{Orientalism examples}: Middle East/East Asia treated as timeless/exotic; stereotypes substitute for actual complexity. \kw{Nader, Studying Up}: anthropologists should study elites, corporations, states, bureaucracies, scientists--not only marginalized ``others.'' \kw{``New savages'' critique}: when anthropologists study West, they may simply relocate othering onto poor/working-class/Western subcultures. \kw{Knowledge ethics}: Is acquiring/sharing knowledge always moral good? For whom? Could knowledge be used for control, military, corporations? \kw{Exam link}: Said + Nader = power shapes both what is known and who becomes object of knowledge.

\section*{W4: CULTURE CONCEPTS}
\kw{Culture is not}: bounded container, homogeneous group essence, static tradition, only non-Western. \kw{Culture is}: contested, learned, symbolic, practiced, power-laden, lived differently by individuals. \kw{Borders}: cultural boundaries are porous; diaspora/globalization complicate nation=culture. \kw{Contestation}: people within same culture disagree. \kw{Individual vs group}: culture shapes choices but does not mechanically determine them. \kw{Carroll}: misunderstandings arise when taken-for-granted assumptions differ; ``home,'' conversation, politeness, directness, privacy vary. \kw{Anthro vs psych}: not just individual motives; asks what shared meanings make behavior sensible.
\sep
\section*{Geertz / Interpretive Anthro}
\kw{Geertz}: culture = webs of significance humans spin; analysis = interpretation of meaning. \kw{Thick description}: describe behavior + cultural meaning + context; wink vs eye twitch. \kw{Thin description}: external act only. \kw{Symbols}: public carriers of meaning; people act through them. \kw{Deep play} (Bentham): stakes so high they seem irrational materially; socially meaningful. \kw{Balinese cockfight}: not ``just gambling''; dramatizes status, masculinity, kinship alliances/rivalries; people bet on roosters owned by kin/allies. \kw{Function}: reinforces social relationships; models hierarchy. \kw{Analogy}: sports fandom/derbies where stakes exceed rational payoff. \kw{Test use}: If question asks why irrational acts persist, answer: symbolic/social value.

\section*{W4-5: KINSHIP BASICS}
\kw{Kinship}: culturally organized relatedness; not reducible to biology. \kw{Schneider}: Western kinship assumed blood+marriage as natural; anthropology must not universalize American assumptions. \kw{``Blood is thicker than water''}: cultural ideology, not universal fact. \kw{Consanguineal}: blood/descent kin. \kw{Affinal/affines}: kin by marriage. \kw{Fictive/chosen kin}: non-biological kin treated as family. \kw{Carsten/adoptees}: kinship made through care, co-residence, memory, documents, searching; blood may matter emotionally but not only factor. \kw{Sahlins, mutuality of being}: kin are people who participate in each other's existence; share experiences, emotions, suffering, food, property, fate. \kw{Hopper/TallBear}: challenge couple/nuclear family norm; love/relations beyond settler sex/family; kinship can be friendship, care networks, poly/chosen relations.
\sep
\section*{Kin Terms / Descent / Marriage}
\kw{Kinship terminology}: languages classify relatives differently; American English vs Tunisian Arabic; names reveal social priorities. \kw{Hawaiian system}: same terms for mother/aunts, father/uncles, siblings/cousins by generation/sex. \kw{Bilateral descent}: both mother/father sides. \kw{Unilateral}: one side. \kw{Patrilineal}: descent through father line. \kw{Matrilineal}: descent through mother line (not same as matriarchy). \kw{Patriarchal}: men hold authority. \kw{Matriarchal}: women hold authority. \kw{Patrilocal/virilocal}: live with husband's family. \kw{Matrilocal}: live with wife's family. \kw{Neolocal}: independent residence. \kw{Nuclear family}: parents+children; not universal default. \kw{OPH}: one-person households. \kw{Marriage}: social contract/covenant/alliance; defines rights/obligations, not just love. \kw{Love vs arranged}: not pure opposites; arranged can involve love, love marriage shaped by social structure. \kw{Endogamy}: marry within group. \kw{Exogamy}: marry outside. \kw{Incest taboo}: widespread but locally defined. \kw{Cross-cousin}: child of parent's opposite-sex sibling. \kw{Parallel-cousin}: child of parent's same-sex sibling. \kw{Sororate}: man marries deceased wife's sister. \kw{Monogamy}: one spouse. \kw{Serial monogamy}: one at a time across life. \kw{Polygamy}: multiple spouses. \kw{Polygyny}: one man + multiple women. \kw{Polyandry}: one woman + multiple men; \kw{fraternal polyandry}: brothers share wife, preserves family property. \kw{Sienna Craig, Ends of Kinship}: Nepali diaspora, marriage practices, kinship across distance; arranged/love/residence patterns adapt transnationally. \kw{Lasar}: many half siblings; family portrait challenges genetic/nuclear assumptions.

\section*{W6: STATE / TRIBE / POWER}
\kw{State vs tribe}: old contrast often oversimplifies; beware evolutionist hierarchy. \kw{Weber}: state claims monopoly on legitimate violence within territory. \kw{Legitimacy}: people accept authority as rightful. \kw{Bureaucracy}: rule by offices, paperwork, procedures; impersonal but staffed by people. \kw{Graeber bureaucracy}: paperwork can produce irrationality, delay, meaningless labor. \kw{State from below} (Holmes): study how policies are lived by migrants/workers; state power felt in bodies, labor, borders, clinics. \kw{Blurry state}: NGOs/corporations/global institutions do state-like work; line between state/economy not obvious. \kw{State as aspirational}: states claim coherence/control they may not achieve.
\sep
\section*{Power Theorists}
\kw{Gramsci, hegemony}: domination through consent/common sense; ruling class ideas become normal. Example: standard language ideology; meritocracy. Hegemony is contestable. \kw{Althusser, ISA}: Ideological State Apparatuses reproduce social order through schools, media, churches, family, unions; people internalize values. \kw{Foucault, sovereign power}: premodern visible punishment/torture/executions display ruler authority. \kw{Discipline}: modern power trains bodies through routines, surveillance, exams, timetables; creates self-regulating subjects. \kw{Panopticon} (Bentham): central tower; prisoners might be watched so behave as if always watched; modern sites: prison, school, hospital, workplace. \kw{Power not only top-down}: diffuse, productive, internalized. \kw{Mitchell, fetishism of state}: common view: state separate entity acting on society. Mitchell: state effect produced through practices/representations. Study: \kw{representations} (what state says it does), \kw{practices} (what officials actually do), \kw{effects} (how people encounter state). Examples: post office, border patrol, documents. \kw{Comparison}: Weber=violence; Gramsci=ideas/consent; Althusser=institutions/ideology; Foucault=discipline/surveillance; Mitchell=state as effect.


\sep
\section*{W1--6 FAST GLOSSARY / APPLICATIONS}
\kw{Social norm example}: silence in library, handshake distance, waiting in line; rule is visible when broken. \kw{Othering example}: describing dentists as ``holy-mouth-men'' or bathrooms as shrines makes Americans seem primitive. \kw{Culture shock}: feeling disoriented when assumptions fail. \kw{Reflexivity}: researcher notices own position/bias; fieldnotes include self as part of data. \kw{Holism}: connect economy, kinship, religion, politics, language, body. \kw{Comparativism}: compare without ranking. \kw{Universalism warning}: do not assume Western family/state/religion/language categories apply everywhere. \kw{Relativism warning}: explain context without excusing harm. \kw{Dialect in Hurston}: evidence of linguistic skill, not ignorance. \kw{Orientalist move}: explain politics as ancient culture rather than occupation/history/power. \kw{Study up examples}: Wall Street, Pentagon, universities, landlords, police, corporations, hospitals. \kw{Deep play answer}: material irrationality becomes socially rational because honor/status/identity are at stake. \kw{Kinship answer}: ask who cares, feeds, lives together, inherits, mourns, sacrifices--not only DNA. \kw{State answer}: ask how state appears in documents, borders, IDs, waiting rooms, offices, police, schools.

\section*{Likely Prompt Skeletons W1--6}
\kw{``Why is Nacirema funny?''} Because it uses exoticizing language on Americans, revealing ethnocentrism and the power of ethnographic description. \kw{``What did Malinowski change?''} From distant/armchair accounts to immersive fieldwork and everyday-life observation. \kw{``How is Hurston both insider/outside?''} Insider by race/language/region familiarity; outsider by academic role/recording; positionality is layered. \kw{``What is Orientalism?''} Not just prejudice; a whole knowledge system that produces the Orient as knowable/inferior and supports rule. \kw{``Culture misunderstanding?''} People act from different background assumptions, so both sides can see themselves as normal/rational. \kw{``Power comparison?''} Weber = legitimate force; Gramsci = consent; Althusser = institutions; Foucault = self-discipline; Mitchell = state effect.

\section*{FRONT QUICK COMPARISONS}
Relativism vs ethnocentrism: understand in context vs judge by own standards. Emic vs etic: insider meaning vs analytic frame. Armchair vs fieldwork: texts from afar vs immersion. Thin vs thick: act alone vs act+meaning. Culture vs biology: learned meanings vs genetic determinism. Hegemony vs discipline: consent/common sense vs self-regulation. Sovereign vs disciplinary power: public punishment vs invisible surveillance. State representation vs practice: claim vs actual operation.


\sep
\section*{W1--6 AUTHOR DETAIL BANK}
\kw{Malinowski, Argonauts}: ethnographer should leave white compound/veranda, live in village, learn language, collect genealogies/maps/cases, record ordinary events as they happen; goal is native's point of view + relation to life. \kw{Boas vs evolutionists}: no single ladder of progress; cultural traits have particular histories; environment matters but does not mechanically determine culture. \kw{Hurston method}: performance/storytelling matters; jokes/folklore are not trivial--they reveal social criticism, pleasure, and identity. \kw{Said method}: read texts/images/scholarship as part of empire; representation produces the object it claims to merely describe. \kw{Nader method}: if power is the problem, research should move upward toward institutions that make decisions. \kw{Carroll examples}: privacy, hospitality, directness, conversation timing, and what counts as ``home'' are culturally organized. \kw{Geertz method}: anthropologist reads culture like text, but interpretation must be grounded in public symbols/practices. \kw{Schneider critique}: ``blood'' and ``law'' are American symbols of kinship, not universal analytic truth. \kw{TallBear critique}: settler colonialism organizes property, sex, marriage, and family; relationality can exceed couple-form. \kw{Holmes}: structural vulnerability; suffering of migrant laborers is produced by hierarchy yet treated as natural/personal.

\section*{W1--6 CONCEPT CHECKLIST}
\kw{Norm}: expected behavior; \kw{sanction}: reaction to violation. \kw{Ritual}: repeated meaningful practice; not only religion. \kw{Magic label}: often used by outsiders to downgrade other practices; compare to Nacirema medicine. \kw{Tradition}: not always ancient/static; can be invented, contested, political. \kw{Local knowledge}: practical and situated; not inferior to science. \kw{Positionality}: race/class/gender/nationality/language affect access and interpretation. \kw{Representation}: who speaks, who is spoken for, which details are selected. \kw{Dialect}: rule-governed language variety; not broken standard language. \kw{Diaspora}: people/relations across homeland and elsewhere; kinship maintained through calls, remittances, visits, marriage. \kw{Settler colonialism}: replacement/dispossession structure, not just past event. \kw{Legibility}: state wants people/forms/land readable through categories, IDs, censuses, maps. \kw{Documentation}: papers make state power real in daily life. \kw{Surveillance}: watching or possibility of watching changes behavior. \kw{Internalization}: power works when people monitor themselves. \kw{Common sense}: hegemony feels natural, not ideological.

\section*{HIGH-YIELD EXAMPLES TO DROP}
\kw{Nacirema}: bathrooms/dentists/hospitals show ethnographic language creates exoticism. \kw{Balinese cockfight}: gambling reveals rank, masculinity, kinship, rivalry. \kw{Standard language}: school teaches elite speech as correct; Gramsci/Althusser link. \kw{Airport/border}: Mitchell state effect + Weber violence + documents. \kw{Library/classroom}: Foucault discipline via silence, seating, grades, surveillance. \kw{Arranged marriage}: challenges Western love-marriage ideology; shows marriage as family/social alliance. \kw{Adoption/half siblings/chosen family}: kinship made through care, recognition, documents, and narrative. \kw{Folklore/jokes}: serious cultural data, not merely entertainment. \kw{Maps/censuses}: state categories create what they claim to measure. \kw{Hospital/clinic}: body becomes site of state/medical authority; compare Nacirema + Holmes.

\section*{MINI CONTRASTS FOR SHORT ANSWERS}
\kw{Relativism vs moral relativism}: method of understanding vs claim that anything is okay. \kw{Culture vs tradition}: culture is broad lived system; tradition is selected/invoked practice. \kw{Kinship vs family}: kinship is analytic system of relatedness; family is one culturally specific form. \kw{Descent vs residence}: who you trace relation through vs where couple lives. \kw{Patrilineal vs patriarchal}: descent through men vs men hold authority. \kw{Matrilineal vs matriarchal}: descent through women vs women hold authority. \kw{Sovereignty vs discipline}: ruler's right to punish vs institutions training conduct. \kw{Representation vs practice}: official image/claim vs everyday action. \kw{Consent vs coercion}: hegemony makes domination feel normal; violence remains backup. \kw{Insider vs native}: insider status is partial; ``native'' label can flatten complexity.

\end{multicols}
\newpage

% ===================== PAGE 2 =====================
\begin{center}\bfseries\fontsize{7}{7}\selectfont CAS AN101 FINAL CHEAT SHEET - BACK: WEEKS 7--12 / GENDER, RELIGION, RACE, ECONOMY, FOOD, MEDICINE, SPACE, SOUND\end{center}\vspace{-4pt}
\begin{multicols*}{4}

\section*{W7: GENDER / SEX / NORMS}
\kw{Sex}: biological traits (but also culturally interpreted). \kw{Gender}: social meanings/roles attached to sexed bodies. \kw{Patriarchy}: male dominance in institutions/family/culture. \kw{Androcentrism}: male as default/standard; women marked as deviation. \kw{Gender binary}: two naturalized categories; media reinforces difference. \kw{Goffman}: ads/images display gender hierarchy (size, touch, gaze, posture, infantilization). \kw{Gendered norms}: expected behavior, emotion, voice, dress, work, sexuality. \kw{Beauty Myth} (Naomi Wolf): beauty standards discipline women; power backlash. \kw{Judith Butler}: gender performative--repeated acts create appearance of stable identity; not fake, socially real through repetition. \kw{Tey Meadow}: transness + ``specter of the regretter''; policy/medicine overfocus on hypothetical regret, shaping access to care. \kw{Ortner}: ``Is female to male as nature is to culture?'' argues many societies symbolically associate women with nature/reproduction and men with culture/public life; asks why hierarchy is constructed. \kw{Martin, Egg and Sperm}: scientific descriptions gender cells: egg passive/damsel, sperm active/heroic; science uses cultural stereotypes. \kw{Blair, Ejaculate Responsibly}: reframes pregnancy responsibility around men/ejaculation; challenges gendered blame. \kw{Putnam/Reeves, Boys Not OK}: masculinity/boys in crisis discourse; analyze gender norms + institutions, not biology alone.
\sep
\section*{Gender Exam Moves}
If image/ad prompt: identify gender display, hierarchy, gaze, activity/passivity. If science/medicine prompt: Martin = biological facts described through gender ideology. If trans prompt: gender norms + institutions regulate bodies. If family/reproduction prompt: link kinship + gender + responsibility.

\section*{W8: RELIGION}
\kw{Religion as category}: contested; not universal/self-evident; Western/Protestant bias often privileges belief, scripture, sincerity, interior faith over ritual/practice/material objects. \kw{Orthodoxy}: claims about correct belief/practice; power decides what counts. \kw{Legal definitions}: courts/states try to define religion, often awkwardly. \kw{Durkheim}: religion divides sacred/profane; social group worships itself through sacred symbols. \kw{Totemic object}: material symbol of clan/society; collective identity. \kw{Pilgrimage}: travel to sacred site; can produce transformation, community, hierarchy, commerce. \kw{Religious liberalism}: modern idea religion as individual choice/interiority; can enable appropriation. \kw{Bucar, Stealing My Religion}: ``playing pilgrim''; borrowing religious practices may erase power/history. \kw{Appropriation vs borrowing}: borrowing = exchange without penalizing donor group; appropriation = host gains status/pleasure while donor group is stigmatized or dispossessed. \kw{Musa, 1452 Poles to Karbala}: ritual/procession/material forms; religion as embodied, political, public. \kw{Key}: religion is lived through bodies, objects, sound, movement, law, power--not just belief.

\section*{W9: RACE / ETHNICITY}
\kw{Race}: socially constructed classification with real effects; not biological essence. \kw{Biological determinism}: false claim race predicts ability/behavior. \kw{Racial commonsense}: everyday ``obvious'' assumptions about race. \kw{Social construct}: made by history/institutions; real consequences. \kw{Reification}: treating constructed category as natural thing. \kw{Passing}: being read/accepted as another race; reveals race depends on perception/social rules. \kw{Hypodescent}: ``one-drop'' rule; mixed ancestry assigned to subordinate race. \kw{Ethnic group}: shared ancestry/culture/history claim. \kw{Ethnicity}: identification/practice, often flexible. \kw{Symbolic ethnicity}: optional/occasional identity practices (food, holidays) especially for privileged groups. \kw{Whiteness}: racial position often unmarked/default; carries privilege. \kw{Garcia}: race/ethnicity are social but embodied in institutions/health/life chances. \kw{NYT Conversation with White People}: whiteness, discomfort, racial ignorance as social position.
\sep
\section*{Falls / Reversal / Slurs}
\kw{Falls, Redneck Customs}: demolition derby as race/class performance; ``redneck'' identity can be reclaimed; ritual/rite of reversal temporarily inverts hierarchy (valuing roughness, destruction, working-class masculinity). \kw{Butler + slurs}: reclaiming resignifies injury, but never fully controlled; context matters. \kw{Bilici, Muslim ethnic comedy}: humor can invert Islamophobia, expose stereotypes, produce solidarity. \kw{Rite of reversal}: temporary inversion of power/social order; not necessarily permanent revolution. Use for comedy, festivals, derby, Key \& Peele names.

\section*{W9: CLASS / BOURDIEU}
\kw{Marxian class}: relation to means of production; bourgeoisie own, proletariat sell labor. \kw{Class consciousness}: awareness of shared class position. \kw{Bourdieu}: class reproduced through capital + habitus + taste. \kw{Economic capital}: money/assets. \kw{Social capital}: networks/connections. \kw{Cultural capital}: valued knowledge, manners, credentials, speech, taste. \kw{Habitus}: embodied dispositions shaped by social position; what feels natural/easy/beautiful. \kw{Taste}: judgments of food/music/art/style mark class; presented as personal preference but socially patterned. \kw{Symbolic violence}: dominated groups accept hierarchy as natural/legitimate. \kw{People Like Us/Tammy's Story}: class visible in speech, home, work, taste; moral judgments attach to poverty. \kw{Language link}: standard language ideology = cultural capital.

\section*{W10: ECONOMY / VALUE}
\kw{Commodity}: object/service produced for exchange. \kw{Use value}: usefulness. \kw{Exchange value}: market value relative to other commodities. \kw{Labor theory of value} (Marx): value rooted in socially necessary labor time. \kw{Surplus value}: profit extracted from labor. \kw{Commodity fetishism}: social relations between workers appear as relations between things; commodity hides labor/exploitation. \kw{Money} (Maurer): not just neutral tool; social, symbolic, moral, emotional. \kw{Fetishization of money}: money treated as if it has independent power. \kw{Capitalism}: wage labor, private ownership, accumulation, markets. \kw{Luxemburg}: capitalism tied to colonial expansion/noncapitalist spaces. \kw{Crisis} (Roitman): ``crisis'' is a narrative framing that implies normality + rupture. \kw{Weber}: Protestant ethic linked religious discipline to capitalist productivity. \kw{Ibn Khaldun}: social cohesion/production/state cycles; early social theory. \kw{Waldfogel, Deadweight Loss of Christmas}: gifts can destroy economic efficiency; contrasts econ rationality with social value. \kw{Adam Grant / productivity}: American moral script equates productivity with goodness. \kw{Cameron Hay-Rollins}: cultural scripts define hard work, worth, deservingness. \kw{Graeber, Bullshit Jobs}: jobs workers believe are pointless; capitalism/bureaucracy produce meaningless labor despite productivity ideals.
\sep
\section*{Mauss / Gift}
\kw{Mauss}: gifts are never free; obligations to give, receive, reciprocate. \kw{Hau}: spirit/force of gift that compels return. \kw{Gift exchange}: creates social bonds, hierarchy, debt, alliance. \kw{Compare}: Marx=commodities hide social relations; Mauss=gifts openly create social relations. \kw{Hidden Brain emotional currency}: money/gifts carry emotion and relationship meaning beyond price.

\section*{W11: FOOD / ANIMALS / ENV}
\kw{Food ideology}: moral/cultural beliefs about eating (clean, natural, ethical, masculine, healthy). \kw{Safran Foer, eating dogs}: disgust is cultural; challenges arbitrary animal categories. \kw{Carol Adams}: ``meat'' erases animal subject; absent referent; links animal suffering and gendered domination. \kw{Cincy Freedom}: escaped cow becomes individual/subject, not meat. \kw{Challenger}: humans as predators have special responsibility due to intelligence/choice. \kw{Berry, disavowal of earth}: modern life hides dependence on land/labor/ecology. \kw{Indigenous relations w/ natural world}: often contrasted w/ Western industrial nature/culture split, but avoid romanticizing. \kw{Anthropocentrism}: human-centered worldview. \kw{Anthropocene}: epoch/idea of massive human planetary impact. \kw{Tsing}: let go of teleology/progress narratives; attend to multispecies worlds, precarity, ruins. \kw{Posthuman imaginary}: decenter human; think animals, plants, objects, environments as actors/relations. \kw{Embodied intelligence}: intelligence not only abstract cognition; bodily/environmental skill. \kw{Rights of Nature film}: legal personhood for rivers/ecosystems; challenges human-only rights.
\sep
\section*{Medical Anthropology}
\kw{Medical anthro}: studies illness/health/healing as biological + cultural + political-economic. Questions: Who gets sick? Who gets care? How is suffering interpreted? How do institutions produce harm? \kw{Holmes, Fresh Fruit}: migrant farmworkers' bodily suffering linked to labor hierarchy, racism, citizenship, clinics; ``naturalized'' inequality. \kw{Intersectionality}: overlapping systems (race/class/gender/citizenship) shape vulnerability. \kw{Symbolic violence}: inequality misrecognized as personal failure/natural difference. \kw{Photo analysis} (Llorens/Sontag): images frame suffering; spectatorship has ethics; photos can reveal or objectify.

\section*{W12: SPACE / PLACE / BODY}
\kw{Yi-Fu Tuan}: space = openness/abstraction; place = space made meaningful through experience, memory, attachment. \kw{Public/private}: culturally made distinction; power decides access, surveillance, respectability. \kw{Space + class}: neighborhoods, homes, schools, stores, transportation encode inequality. \kw{Habitus + techniques of body}: learned ways of sitting, walking, speaking, eating, studying; class/culture embodied. \kw{Home} (Carroll): not universal meaning; privacy, hospitality, comfort differ. \kw{Conversation} (Carroll): turn-taking, silence, directness, interruption vary cross-culturally.
\sep
\section*{Time / History}
\kw{Time}: spoken as space (ahead/behind, long/short, moving forward); socially organized by calendars, schedules, waiting, deadlines. \kw{Jackson, temporal experience}: people live time differently; past/present/future shaped by memory, suffering, hope, politics. \kw{Bringing past into present}: monuments, rituals, museums, stories, land claims, law, trauma, holidays. \kw{History writing}: not neutral; selects perspective. \kw{Silverman, This Land is Their Land}: critical history of U.S. founding/Indigenous dispossession; challenges celebratory national narratives. \kw{Critical history}: asks whose past is centered, erased, or justified.

\section*{Sound / Silence}
\kw{John Cage, 4'33''}: silence is never empty; listening reveals ambient sound and expectations about music/art. \kw{Basso}: silence has social meaning; not merely absence; can show respect, anger, uncertainty depending context. \kw{Sacred sounds}: sound can create religious authority, atmosphere, belonging. \kw{Goshadze, Ghana harvest festival}: ``quiet'' vs ``not-noise''; sound judged by social meaning, not just decibels. \kw{Sound claims space}: music, bells, calls, chants, noise complaints mark belonging/power. \kw{Sonic experience}: hearing is cultural/political.

\section*{LANGUAGE / POWER (FINAL LECTURES)}
\kw{Standard language ideology}: belief there is one correct way to speak; privileges elite speech and marginalizes dialects. \kw{Language stratification}: varieties ranked by power, not logic. \kw{Cultural capital}: knowing prestigious pronunciation/names/language gives access. \kw{Key \& Peele substitute teacher}: joke is reversal; white names revealed as convention too; de-centers white norms. \kw{Mock Spanish} (Jane Hill): ``Cinco de Drinko,'' adding -o; allows white speakers to appear relaxed/funny while Spanish speakers are stereotyped as lazy/party/drunk. \kw{Appropriation vs borrowing}: host group gains value; donor group punished for same practice. \kw{AAE example}: used by outsiders for humor/coolness/sexual sophistication while Black speakers face discrimination. \kw{LGBTQ language}: ``yas queen'' can signal safe/cosmopolitan identity for outsiders; context matters. \kw{Young, Should Writers Use They Own English?}: challenges standard English dominance; code-meshing; language justice.
\sep
\section*{Military Anthropology / Ethics}
\kw{Human Terrain System}: social scientists embedded with military in Iraq/Afghanistan. Debate: applied anthro to reduce harm vs complicity in occupation/war. \kw{Gusterson}: problem not lack of cultural understanding but occupation; anthropologists should study culture of national security state. \kw{AAA}: broad consensus against anthropologists as military tools. \kw{Big Q}: What is knowledge for? For whom? Does ``getting hands dirty'' help or legitimate violence?


\sep
\section*{W7--12 FAST GLOSSARY / APPLICATIONS}
\kw{Androcentrism example}: ``mankind,'' male crash-test dummy, male symptoms as medical default. \kw{Performative example}: clothes, voice, posture, pronouns, repeated recognition make gender feel natural. \kw{Religion prompt}: ask who has authority to define sacred/orthodox/real religion. \kw{Appropriation prompt}: ask whether donor group is respected/paid/free to practice or stigmatized. \kw{Race prompt}: explain how category is constructed, then show real institutional effects. \kw{Class prompt}: taste/language/credentials feel personal but reproduce hierarchy. \kw{Commodity prompt}: identify hidden labor/social relations behind object. \kw{Gift prompt}: identify obligation/debt/status, not price. \kw{Food prompt}: disgust and ethics reveal culture; animal categories are not natural. \kw{Medical prompt}: connect body suffering to labor, policy, race/class/citizenship. \kw{Space/place prompt}: ask how memory/power turns location into meaningful place. \kw{Sound prompt}: not all sound is ``noise''; judgment depends on belonging/power/context. \kw{Photo prompt}: who is pictured, by whom, for whom, what emotion/politics produced?

\section*{More Author Hooks}
\kw{Ehrenreich, Nickel and Dimed}: low-wage work is exhausting and expensive; poverty is produced, not moral failure. \kw{Bourgois/Schonberg}: addiction/homelessness through structural violence, not individual pathology alone. \kw{Meneley}: sociability and hierarchy in Yemeni town; value created through hosting, exchange, reputation. \kw{Abu-Lughod}: Bedouin poetry expresses sentiment indirectly where honor/modesty restrict open emotion; language reveals gender/power. \kw{Invisibilia emotions}: emotions are culturally learned/managed, not purely internal biology. \kw{Hidden Brain money}: exchange can damage or express relationships depending framing. \kw{4'33''}: listening practice reveals social expectations around art/noise/silence. \kw{Human Terrain}: ethics of applied knowledge under military occupation.

\section*{READING INDEX - ONE-LINE TRIGGERS}
Miner=Nacirema/ethnocentrism. Boas=relativism/anti-racism/Eskimo. Malinowski=fieldwork/imponderabilia. Hurston=dialect/insider/burden. Said=Orientalism/power-knowledge. Nader=study up. Carroll=misunderstandings/home/conversation. Geertz=thick description/deep play. Schneider=kinship not biology. Lasar=half siblings/nontraditional family. Hopper=chosen love/against couple norm. TallBear=relations beyond settler family. Craig=diaspora kinship. Foucault=punishment/panopticon. Holmes=migrant suffering/state below. Ortner=female/nature. Putnam/Reeves=boys crisis. Blair=male responsibility. Martin=gendered science. Musa=Karbala ritual. Durkheim=totem/social cohesion. Bucar=religious appropriation. Garcia=race/ethnicity. Falls=redneck/reversal. Bourdieu=capital/habitus/taste. Marx=commodity/labor/fetishism. Mauss=gift/hau/obligation. Graeber=bullshit jobs/bureaucracy. Ehrenreich=low-wage labor. Adams=meat/absent referent. Bourgois/Schonberg=addiction/poverty/structural violence. Meneley=value/hierarchy/sociability. Jackson=time. Silverman=critical Indigenous history. Goshadze=quiet/not-noise. Bilici=Muslim comedy. Hill=mock Spanish. Young=language justice. Abu-Lughod=Bedouin poetry/sentiment; indirect expression reveals honor/gender/social hierarchy.

\section*{EXAM ANSWER FORMULAS}
\kw{Define}: term + theorist + contrast + example. \kw{Apply}: identify case details, name concept, show how concept explains meaning/power. \kw{Compare}: same problem, different mechanism. \kw{Good pairs}: Boas vs ethnocentrism; Geertz vs thin description; Weber/Gramsci/Foucault/Mitchell; Marx vs Mauss; borrowing vs appropriation; race as construct vs biological determinism; space vs place; sex vs gender. \kw{Universal warning}: avoid saying ``they believe weird things''; write ``within this system, practice means X and does Y.'' \kw{Power warning}: ask who benefits, who is represented, who is silenced, what institution reproduces it. \kw{Culture warning}: culture is not an excuse; it is context, meaning, practice, and power.

\end{multicols*}
\end{document}
